\documentclass[10pt,a4paper]{amsart}
\usepackage[margin=19mm]{geometry}
\usepackage{amsmath,amssymb,mathtools}
\usepackage[scr=rsfs,cal=euler]{mathalfa}
\usepackage{xfrac}
\usepackage[unicode,hidelinks,bookmarks=false]{hyperref}
\newcommand{\ie}{i.e.\ }
\newcommand{\cf}{cf.\ }
\newcommand{\resp}{respectively\ }
\newcommand{\Subsection}{\subsection}
\providecommand{\nobreakdash}{\nobreak\hbox{-}}
\setlength{\emergencystretch}{2em}
\multlinegap0pt
\usepackage{bm}
\usepackage{bbm}
\usepackage{dsfont}
\usepackage{xspace}
\usepackage{cancel}
\usepackage{mathtools}
\usepackage{amsthm}
\makeatletter
\@ifundefined{theorem}{\newtheorem{theorem}{Theorem}}{}
\@ifundefined{proposition}{\newtheorem{proposition}[theorem]{Proposition}}{}
\makeatother
\makeatletter
\newcommand{\MM}{\mathbf{M}}
\expandafter\ifx\csname sageexplain\endcsname\relax
\newenvironment{sageexplain}
{
  \vskip 1em
  \begin{addmargin}[0.05\linewidth]{0.05\linewidth}
  \footnotesize
  \color{black!70}
  \parindent=0pt
  \refstepcounter{sageexplaincounter}%
  \textbf{SageMath example~\thesageexplaincounter.}
}
{
  \end{addmargin}
  \vskip 1em
}
\fi
\makeatother
\title[A computational approach to Drinfeld modules]{A computational approach to Drinfeld modules}
\author{C\'ecile Armana, Elena Berardini, Xavier Caruso, Antoine Leudi\`ere, Jade Nardi, Fabien Pazuki}
\date{Source: arXiv:2601.02162v1 (CC BY 4.0)}
\begin{document}
\maketitle

\noindent\emph{Source and licence.} A computational approach to Drinfeld modules. Version arXiv:2601.02162v1 (CC BY 4.0), distributed under the Creative Commons Attribution 4.0 International licence, as recorded in the arXiv licence metadata for this exact version and shown on the abstract page https://arxiv.org/abs/2601.02162v1. Full rights notice: licenses/arxiv-2601.02162-CC-BY-4.0.txt.\par\smallskip

\noindent\textit{Editorial scope.} Counted unit: the Proposition whose source label is \texttt{prop:basismotive} in the article (source0.tex lines 1467--1472, source label \texttt{prop:basismotive}), delivered together with the article's own complete original proof of it (source0.tex lines 1474--1497, one \texttt{proof} environment of 24 lines). No mathematics is written, repaired or re-derived: statement and proof are the article's own text line for line. Required context retained uncounted: none. Every definition, display and label of those lines is retained unchanged. Omitted from this package and not redistributed: 1--1466, 1473--1473, 1498--4431. Nothing inside a retained excerpt is omitted or altered: every retained body is delivered exactly as the article's lines are written, so no omission receipt of any kind accompanies this package. Citations: the retained excerpts carry no citation command at all, so no bibliography entry of the article is transcribed and no reference is redistributed. Reference adaptation: none was needed and none was made; every reference inside the delivered unit resolves inside it, so no unresolved reference or citation is produced. Printed slips kept literal: no printed word, symbol, hypothesis or number of the counted statement or of its proof was altered. Layout and class: the article's own class is \texttt{article} with packages including fontenc, inputenc, mathpazo, geometry, amsmath, amsfonts, amssymb, amsthm, mathtools, soul, hyperref, tikz, framed, enumitem; the wrapper is the lane's pinned \texttt{amsart} editorial wrapper at 10pt on A4, which restates the theorem-like environments the retained text needs and adds the article's own macro definitions that the retained text needs (\texttt{\textbackslash MM}), with extra wrapper packages (bm, bbm, dsfont, xspace, cancel, mathtools). The delivered numbering is the unit's own. Assets: no retained span contains a figure environment, a graphics-inclusion command, a PDF-page inclusion, a table or a TikZ picture; no image file is copied into this package, the package declares no asset dependency and no image file is redistributed with it. Licence: version arXiv:2601.02162v1 of this article is distributed under the Creative Commons Attribution 4.0 International licence, as recorded in the arXiv licence metadata for this exact version and shown on the abstract page https://arxiv.org/abs/2601.02162v1. A full rights notice is delivered at licenses/arxiv-2601.02162-CC-BY-4.0.txt. Characters: every retained excerpt is pure ASCII, so the delivered engine, XeLaTeX with its default Latin Modern text font, reports no missing character.
\begin{proposition}
\label{prop:basismotive}
If $\phi : A \to F\{\tau\}$ is a Drinfeld module of rank $r$,
then $\MM(\phi)$ is free of rank $r$ over $F[T]$ with basis
$(1, \tau, \ldots, \tau^{r-1})$.
\end{proposition}

\begin{proof}
We explain how to decompose an element $m \in \MM(\phi)$ in the
given basis. We recall that, by definition, $m$ is an Ore polynomial
in $F\{\tau\}$.
The latter being right Euclidean, we can decompose $m$ in the basis $\phi_T$ as follows:
$$m = m_0 + m_1 \phi_T + \cdots + m_k \phi_T^k,$$
for some nonnegative integer $k$ and some uniquely determined
$m_0, \ldots, m_k \in F\{t\}$ of degree at most $r{-}1$ where
$r = \deg_\tau \phi_T$ is the rank of $\phi$.
Indeed, $m_0$ is obtained as the remainder in the right division of
$m$ by $\phi_T$ and the next $m_i$ are obtained similarly by
dividing the successive quotients. Now writing
$$m_i =
\lambda_{i,0} + \lambda_{i,1} \tau + \cdots
\lambda_{i,r-1} \tau^{r-1},$$ we get the formula:
$$m = \sum_{i=0}^k \sum_{j=0}^{r-1} \lambda_{i,j} \tau^j \phi_T^i
    = \sum_{j=0}^{r-1} \left(\sum_{i=0}^k \lambda_{i,j} T^i\right) \tau^j,$$
which turns out to be the decomposition of $m$ in the basis
$(1, \tau, \ldots, \tau^{r-1})$.

The previous calculation shows that the family $(1, \tau, \ldots,
\tau^{r-1})$ generates $\MM(\phi)$. Proving that it is free is left
as an easy exercise for the reader.
\end{proof}

\end{document}
